% 子文件：5.导数与微分.tex
% 对应教材（同济《高等数学》第七版·上册）第二章 导数与微分
%   第一节 导数概念 / 第二节 函数的求导法则 / 第三节 高阶导数
%   第四节 隐函数及由参数方程所确定的函数的导数·相关变化率 / 第五节 函数的微分
% 主文档使用 % 子文件：5.导数与微分.tex
% 对应教材（同济《高等数学》第七版·上册）第二章 导数与微分
%   第一节 导数概念 / 第二节 函数的求导法则 / 第三节 高阶导数
%   第四节 隐函数及由参数方程所确定的函数的导数·相关变化率 / 第五节 函数的微分
% 主文档使用 % 子文件：5.导数与微分.tex
% 对应教材（同济《高等数学》第七版·上册）第二章 导数与微分
%   第一节 导数概念 / 第二节 函数的求导法则 / 第三节 高阶导数
%   第四节 隐函数及由参数方程所确定的函数的导数·相关变化率 / 第五节 函数的微分
% 主文档使用 % 子文件：5.导数与微分.tex
% 对应教材（同济《高等数学》第七版·上册）第二章 导数与微分
%   第一节 导数概念 / 第二节 函数的求导法则 / 第三节 高阶导数
%   第四节 隐函数及由参数方程所确定的函数的导数·相关变化率 / 第五节 函数的微分
% 主文档使用 \input{5.导数与微分.tex} 引入，不含 \documentclass / \begin{document}

\section{导数概念}

\subsection{两个引例}
引例一（变速直线运动的速度）：位移 $s=s(t)$，从 $t_0$ 到 $t_0+\Delta t$ 的平均速度 $\bar v=\dfrac{\Delta s}{\Delta t}$，令 $\Delta t\to0$ 得瞬时速度 $v(t_0)=\lim\limits_{\Delta t\to0}\dfrac{\Delta s}{\Delta t}$。

引例二（平面曲线的切线斜率）：曲线 $y=f(x)$ 上点 $M(x_0,y_0)$ 处切线斜率 $k=\lim\limits_{\Delta x\to0}\dfrac{\Delta y}{\Delta x}$。

两个引例的形式完全相同：都是\textbf{函数增量与自变量增量之比当自变量增量趋于零时的极限}。抽去具体意义就得到导数概念。

\subsection{导数的定义}
\textbf{定义}\quad 设函数 $y=f(x)$ 在点 $x_0$ 的某个邻域内有定义，当自变量 $x$ 在 $x_0$ 处取得增量 $\Delta x$（$x_0+\Delta x$ 仍在该邻域内）时，相应地函数取得增量
\[
\Delta y=f(x_0+\Delta x)-f(x_0);
\]
如果 $\Delta y$ 与 $\Delta x$ 之比当 $\Delta x\to0$ 时的极限存在，则称函数 $y=f(x)$ 在点 $x_0$ 处\textbf{可导}，并称这个极限为函数 $y=f(x)$ 在点 $x_0$ 处的\textbf{导数}，记为 $y'|_{x=x_0}$、$\dfrac{\mathrm{d}y}{\mathrm{d}x}\Big|_{x=x_0}$ 或 $f'(x_0)$，即
\keyformula{f'(x_0)=\lim_{\Delta x\to0}\frac{\Delta y}{\Delta x}=\lim_{\Delta x\to0}\frac{f(x_0+\Delta x)-f(x_0)}{\Delta x}}
\noindent 也可写成等价形式（令 $x=x_0+\Delta x$）
\keyformula{f'(x_0)=\lim_{x\to x_0}\frac{f(x)-f(x_0)}{x-x_0}=\lim_{h\to0}\frac{f(x_0+h)-f(x_0)}{h}}
\noindent 若上述极限不存在，则称 $f(x)$ 在 $x_0$ 处不可导；若极限为无穷大，习惯上也称导数不存在（导数为无穷大）。

\textbf{导函数}\quad 若 $f(x)$ 在区间 $I$ 内每一点都可导，则对每个 $x\in I$ 都有唯一确定的导数值 $f'(x)$ 与之对应，由此得到 $I$ 上的新函数，称为 $f(x)$ 的导函数，记作 $y'$、$f'(x)$ 或 $\dfrac{\mathrm{d}y}{\mathrm{d}x}$。

\subsubsection{单侧导数}
\[\text{左导数：}f'_-(x_0)=\lim_{\Delta x\to0^-}\frac{f(x_0+\Delta x)-f(x_0)}{\Delta x};\qquad
\text{右导数：}f'_+(x_0)=\lim_{\Delta x\to0^+}\frac{f(x_0+\Delta x)-f(x_0)}{\Delta x}.\]
\keyformula{f'(x_0)\ \text{存在}\iff f'_-(x_0)=f'_+(x_0)\ \text{均存在且相等}}

\corepoint{求分段函数在\textbf{分界点}处的导数，必须用定义分左右两侧算 $f'_-$、$f'_+$ 再比较，不能用求导公式直接代——分界点两侧的表达式不同。而在各段\textbf{内部}可用公式求导。$f$ 在 $[a,b]$ 上可导指：在 $(a,b)$ 内可导，且 $f'_+(a)$、$f'_-(b)$ 存在。}

\subsection{导数的几何意义与物理意义}
$f'(x_0)$ 是曲线 $y=f(x)$ 在点 $M(x_0,y_0)$ 处\textbf{切线的斜率}。于是
\keyformula{\text{切线：}\ y-y_0=f'(x_0)(x-x_0);\qquad \text{法线：}\ y-y_0=-\frac{1}{f'(x_0)}(x-x_0)\ (f'(x_0)\neq0)}
\noindent 若 $f'(x_0)=\infty$，则切线为竖直直线 $x=x_0$。

一般地，$f'(x_0)$ 就是因变量对自变量的变化率（瞬时变化率、边际）：$v=\frac{\mathrm{d}s}{\mathrm{d}t}$ 是速度，$\rho=\frac{\mathrm{d}Q}{\mathrm{d}p}$ 是边际需求等。

\subsection{函数可导性与连续性的关系}
\textbf{定理}\quad 若函数 $y=f(x)$ 在点 $x_0$ 处可导，则它在点 $x_0$ 处必连续。

\corepoint{逆命题不成立：连续不一定可导。典型反例 $y=|x|$ 在 $x=0$ 处连续，但 $f'_-(0)=-1\neq f'_+(0)=1$，不可导。$\sqrt[3]{x}$ 在 $x=0$ 处连续且 $f'(0)=\infty$（有竖直切线）。“可导 $\Rightarrow$ 连续”常用其逆否形式：\textbf{不连续必不可导}。判断“在 $x_0$ 是否可导”的标准流程是：先看连续（不连续直接不可导）→ 再用定义算左右导数比较。}

\section{函数的求导法则}

\subsection{函数的和、差、积、商的求导法则}
设 $u=u(x)$、$v=v(x)$ 都在点 $x$ 处可导，则
\keyformula{(u\pm v)'=u'\pm v',\qquad (uv)'=u'v+uv',\qquad \Bigl(\frac uv\Bigr)'=\frac{u'v-uv'}{v^2}\ (v\neq0)}
\noindent 推论：$(Cu)'=Cu'$；$\Bigl(\dfrac{1}{v}\Bigr)'=-\dfrac{v'}{v^2}$。

\corepoint{乘积法则不能写成 $u'v'$；三个因子相乘要两项三项轮流求导：$(uvw)'=u'vw+uv'w+uvw'$。商的分子顺序是 $u'v-uv'$，与积的区别就是“减号在前”，记错会使答案差一个负号。}

\subsection{反函数的求导法则}
\textbf{定理}\quad 若函数 $x=\varphi(y)$ 在某区间内单调、可导且 $\varphi'(y)\neq0$，则它的反函数 $y=f(x)$ 在该对应区间内也可导，且
\keyformula{f'(x)=\frac{1}{\varphi'(y)}\qquad\text{即}\qquad \frac{\mathrm{d}y}{\mathrm{d}x}=\frac{1}{\frac{\mathrm{d}x}{\mathrm{d}y}}}

\corepoint{反函数的导数等于原函数导数的倒数，但两边要“换变量”：算完 $\frac{1}{\varphi'(y)}$ 后必须把 $y$ 用 $y=f(x)$ 回代成 $x$ 的表达式。这个法则主要用于\textbf{推导}反三角函数、对数函数的导数公式（因为它们的原函数好求导），而不是用来做具体计算。}

\subsection{复合函数的求导法则（链式法则）}
\textbf{定理}\quad 若 $u=g(x)$ 在点 $x$ 处可导，$y=f(u)$ 在对应点 $u=g(x)$ 处可导，则复合函数 $y=f[g(x)]$ 在点 $x$ 处可导，且
\keyformula{\frac{\mathrm{d}y}{\mathrm{d}x}=f'(u)\cdot g'(x)\qquad\text{或}\qquad y'_x=y'_u\cdot u'_x}
\noindent 多层复合逐层相乘：$y=f(u)$，$u=g(v)$，$v=h(x)$，则 $y'_x=f'(u)\,g'(v)\,h'(x)$。

\corepoint{链式法则三条铁律：(1) 不要跳层，有几层就乘几个因子；(2) 每层都“对外层变量求导，再乘内层对自变量的导数”；(3) 用 $y'_x=y'_u u'_x$ 写法时，最后要把 $u$ 回代成 $x$ 的表达式。$(\sin^2x)'=2\sin x\cos x$ 不是 $2\sin x$；$(e^{x^2})'=2xe^{x^2}$ 不要漏掉 $2x$。}

\subsection{基本求导法则与导数公式}
\begin{center}
\begin{tabular}{@{}ll@{\hspace{2em}}ll@{}}
\toprule
$f(x)$ & $f'(x)$ & $f(x)$ & $f'(x)$ \\
\midrule
$C$ & $0$ & $\sin x$ & $\cos x$ \\
$x^{\mu}$ & $\mu x^{\mu-1}$ & $\cos x$ & $-\sin x$ \\
$a^{x}$ & $a^{x}\ln a$ & $\tan x$ & $\sec^2x$ \\
$e^{x}$ & $e^{x}$ & $\cot x$ & $-\csc^2x$ \\
$\log_a x$ & $\dfrac{1}{x\ln a}$ & $\sec x$ & $\sec x\tan x$ \\
$\ln x$ & $\dfrac1x$ & $\csc x$ & $-\csc x\cot x$ \\
$\arcsin x$ & $\dfrac{1}{\sqrt{1-x^2}}$ & $\arctan x$ & $\dfrac{1}{1+x^2}$ \\
$\arccos x$ & $-\dfrac{1}{\sqrt{1-x^2}}$ & $\operatorname{arccot}x$ & $-\dfrac{1}{1+x^2}$ \\
\bottomrule
\end{tabular}
\end{center}
\noindent 熟记这张表是求一切导数的基础；配合四则法则、链式法则、反函数法则，就可求初等函数在定义区间内的任意阶导数。

\section{高阶导数}

\subsection{定义与记法}
函数 $y=f(x)$ 的导数 $f'(x)$ 仍是 $x$ 的函数。若 $f'(x)$ 在点 $x$ 处可导，则称 $f'(x)$ 的导数为 $f(x)$ 的\textbf{二阶导数}，记作
\[
y'',\quad f''(x),\quad \frac{\mathrm{d}^2y}{\mathrm{d}x^2};
\]
类似地定义三阶、四阶导数，$n$ 阶导数记作 $y^{(n)}$、$f^{(n)}(x)$、$\dfrac{\mathrm{d}^ny}{\mathrm{d}x^n}$。二阶及二阶以上的导数统称\textbf{高阶导数}。

\subsection{几个常用的 $n$ 阶导数公式}
\keyformula{\begin{aligned}
(e^{ax})^{(n)}&=a^{n}e^{ax}, &\quad (\sin x)^{(n)}&=\sin\Bigl(x+\frac{n\pi}{2}\Bigr),\\
(\cos x)^{(n)}&=\cos\Bigl(x+\frac{n\pi}{2}\Bigr), &\quad (x^{\mu})^{(n)}&=\mu(\mu-1)\cdots(\mu-n+1)x^{\mu-n},\\
(\ln x)^{(n)}&=\frac{(-1)^{n-1}(n-1)!}{x^{n}}. & &
\end{aligned}}

\corepoint{求 $n$ 阶导数的一般方法是“算前几阶、找规律、用数学归纳法证明”。$\sin x,\cos x$ 每求导一次相位前进 $\frac\pi2$，这是周期性函数的高阶导数的常用技巧；$n$ 从 $1$ 开始时 $\sin^{(n)}$ 的公式同样成立。}

\subsection{莱布尼茨公式}
若 $u=u(x)$、$v=v(x)$ 均有 $n$ 阶导数，则
\keyformula{(uv)^{(n)}=\sum_{k=0}^{n}\mathrm{C}_n^{k}\,u^{(n-k)}\,v^{(k)}}
\noindent 其中 $u^{(0)}=u$，$v^{(0)}=v$。形式与二项式展开定理完全相同。

\corepoint{莱布尼茨公式适用于“两个因子之积、其中一个因子（如 $x^m$、多项式）求若干阶导后变零”的情形：多项式 $x^m$ 的 $n$ 阶导在 $n>m$ 时为零，所以和式只剩前面有限项。用错方向的典型是把 $e^x\ln(1+x)$ 硬套莱布尼茨——那不划算。}

\section{隐函数及由参数方程所确定的函数的导数·相关变化率}

\subsection{隐函数的导数}
若方程 $F(x,y)=0$ 确定了 $y$ 是 $x$ 的函数，则称 $y$ 为 $x$ 的\textbf{隐函数}。求导方法：把 $y$ 看作 $x$ 的函数（$y=y(x)$），对方程两边同时关于 $x$ 求导，遇到 $y$ 时按链式法则乘以 $y'$，得到含 $y'$ 的方程，解出 $y'$。

\corepoint{隐函数求导的三个要点：(1) 记住 $y$ 是 $x$ 的函数，$\frac{\mathrm{d}}{\mathrm{d}x}[f(y)]=f'(y)y'$，如 $\frac{\mathrm{d}}{\mathrm{d}x}(y^2)=2yy'$；(2) 结果是含 $x,y$ 的式子，通常保留 $y$（不必解出 $y$）；(3) 求二阶导时，把一阶导的结果\textbf{代入}后再求导，不要停下。求曲线 $F(x,y)=0$ 在点 $(x_0,y_0)$ 处的切线，可用 $y'|_{(x_0,y_0)}$ 直接代入。}

\subsubsection{对数求导法}
对 $y=u(x)^{v(x)}$（$u>0$）或幂连乘、连除式，先在两边取对数 $\ln y=v\ln u$，再两边求导：
\[
\frac{y'}{y}=v'\ln u+v\frac{u'}{u}\ \Longrightarrow\ y'=u^{v}\Bigl(v'\ln u+\frac{vu'}{u}\Bigr).
\]

\corepoint{对数求导法的适用信号：幂指函数（底与指数都含 $x$）、多个因子连乘连除、含根式与乘方。它把乘除化为加减、把指数化为乘数。结果是 $y'$ 的表达式里要\textbf{回乘 $y$ 本身}，别漏掉这个因子（忘记回代是最常见错误）。}

\subsection{由参数方程所确定的函数的导数}
设 $x=\varphi(t)$、$y=\psi(t)$ 都可导且 $\varphi'(t)\neq0$，则
\keyformula{\frac{\mathrm{d}y}{\mathrm{d}x}=\frac{\frac{\mathrm{d}y}{\mathrm{d}t}}{\frac{\mathrm{d}x}{\mathrm{d}t}}=\frac{\psi'(t)}{\varphi'(t)}}
\noindent 二阶导数（对 $x$ 再求一次导，注意 $x$ 仍是 $t$ 的函数）：
\keyformula{\frac{\mathrm{d}^2y}{\mathrm{d}x^2}=\frac{\mathrm{d}}{\mathrm{d}t}\Bigl(\frac{\mathrm{d}y}{\mathrm{d}x}\Bigr)\cdot\frac{1}{\varphi'(t)}
=\frac{\psi''(t)\varphi'(t)-\psi'(t)\varphi''(t)}{\bigl[\varphi'(t)\bigr]^{3}}}

\corepoint{二阶导数\textbf{绝不能}写成 $\dfrac{\psi''(t)}{\varphi''(t)}$。正确做法：先把一阶导 $\dfrac{\mathrm{d}y}{\mathrm{d}x}$ 看成 $t$ 的函数，对它关于 $t$ 求导，再乘以 $\dfrac{\mathrm{d}t}{\mathrm{d}x}=\dfrac{1}{\varphi'(t)}$。这一条几乎每份卷子都在考。}

\subsection{相关变化率}
设 $x=x(t)$、$y=y(t)$ 都可导，而 $x$ 与 $y$ 之间存在关系 $F(x,y)=0$。对 $t$ 求导得
\[
\frac{\mathrm{d}F}{\mathrm{d}x}\cdot\frac{\mathrm{d}x}{\mathrm{d}t}+\frac{\mathrm{d}F}{\mathrm{d}y}\cdot\frac{\mathrm{d}y}{\mathrm{d}t}=0,
\]
从而可由一个变化率求出另一个变化率。这类问题称为\textbf{相关变化率}问题。

\corepoint{解相关变化率问题的步骤：(1) 建立变量之间的方程（几何关系、物理定律）；(2) 两边同时对 $t$ 求导（所有变量都是 $t$ 的函数）；(3) 代入某一时刻的已知值求出待求变化率。注意“代入数值”必须放在求导\textbf{之后}——先代入会把变量变成常数而使导数为 0。}

\section{函数的微分}

\subsection{微分的定义}
\textbf{定义}\quad 设函数 $y=f(x)$ 在某区间内有定义，$x_0$ 及 $x_0+\Delta x$ 都在该区间内。如果函数的增量可表示为
\[
\Delta y=f(x_0+\Delta x)-f(x_0)=A\,\Delta x+o(\Delta x),
\]
其中 $A$ 是不依赖于 $\Delta x$ 的常数，则称 $f(x)$ 在点 $x_0$ 处\textbf{可微}，而 $A\,\Delta x$ 称为 $f(x)$ 在点 $x_0$ 处的\textbf{微分}，记作 $\mathrm{d}y$，即 $\mathrm{d}y=A\,\Delta x$。

\keyformula{f(x)\ \text{在}\ x_0\ \text{处可微}\iff f(x)\ \text{在}\ x_0\ \text{处可导，且}\ \mathrm{d}y=f'(x_0)\,\Delta x}

\noindent 又因 $\mathrm{d}x=\Delta x$（自变量的微分等于自变量的增量），所以
\keyformula{\mathrm{d}y=f'(x)\,\mathrm{d}x,\qquad \frac{\mathrm{d}y}{\mathrm{d}x}=f'(x)}
\noindent 后一式说明：\textbf{导数就是函数的微分与自变量的微分之商}，这就是“微商”一词的来历。

\corepoint{一元函数“可微 $\iff$ 可导”，二者是同一件事；但 $\Delta y$ 与 $\mathrm{d}y$ 不同：$\Delta y$ 是真正的函数增量，$\mathrm{d}y=f'(x)\Delta x$ 是它的\textbf{线性主部}，$\Delta y=\mathrm{d}y+o(\Delta x)$。$\mathrm{d}y$ 依赖于 $x$ 与 $\Delta x$ 两个量，$\mathrm{d}x$ 与 $\Delta x$ 恒等。求 $\mathrm{d}y$ 时不要忘记最后的 $\mathrm{d}x$。}

\subsection{微分的几何意义}
在曲线 $y=f(x)$ 上取点 $M(x_0,y_0)$，当横坐标增加 $\Delta x$ 时，曲线上的纵坐标增量为 $\Delta y=NT$，而切线上相应的纵坐标增量为 $\mathrm{d}y$。因此 $\mathrm{d}y$ 是曲线在点 $M$ 处\textbf{切线的纵坐标增量}，即“以直代曲”的线性近似。

\subsection{基本初等函数的微分公式与微分运算法则}
把求导公式两边同乘 $\mathrm{d}x$ 即得微分公式，如
\[
\mathrm{d}(x^{\mu})=\mu x^{\mu-1}\mathrm{d}x,\quad \mathrm{d}(\sin x)=\cos x\,\mathrm{d}x,\quad
\mathrm{d}(\arctan x)=\frac{\mathrm{d}x}{1+x^2},\quad \mathrm{d}(\ln x)=\frac{\mathrm{d}x}{x}.
\]
微分运算法则：
\keyformula{\mathrm{d}(u\pm v)=\mathrm{d}u\pm\mathrm{d}v,\qquad \mathrm{d}(uv)=v\,\mathrm{d}u+u\,\mathrm{d}v,\qquad \mathrm{d}\Bigl(\frac uv\Bigr)=\frac{v\,\mathrm{d}u-u\,\mathrm{d}v}{v^{2}}}

\subsubsection{复合函数的微分与一阶微分形式不变性}
设 $y=f(u)$ 可微，$u=g(x)$ 可微，则 $\mathrm{d}y=f'(u)\,\mathrm{d}u$。注意这里 $u$ 既可以是自变量，也可以是另一个变量的函数，形式都是 $f'(u)\mathrm{d}u$，这就是\textbf{一阶微分形式不变性}。

\corepoint{一阶微分形式不变性指的是“形式”不变（$\mathrm{d}y=f'(u)\mathrm{d}u$ 恒成立），不是说导数不变（$y'_x=y'_u u'_x$ 差一个 $u'_x$）。它使凑微分（如 $\cos2x\,\mathrm{d}x=\frac12\mathrm{d}(\sin2x)$）成为可能，是第 4 章换元积分法的基础。}

\subsection{微分在近似计算中的应用}
当 $|\Delta x|$ 很小时，
\keyformula{\Delta y\approx \mathrm{d}y=f'(x_0)\Delta x,\qquad f(x_0+\Delta x)\approx f(x_0)+f'(x_0)\Delta x}
\noindent 特别地，取 $x_0=0$ 得 $f(x)\approx f(0)+f'(0)x$，由此得到常用近似式：
\keyformula{(1+x)^{\alpha}\approx1+\alpha x,\quad \sin x\approx x,\quad \tan x\approx x,\quad e^{x}\approx1+x,\quad \ln(1+x)\approx x\quad(|x|\ \text{很小})}

\noindent 误差估计：绝对误差 $|\Delta y|\approx|\mathrm{d}y|=|f'(x)|\,\Delta x$，相对误差 $\dfrac{|\mathrm{d}y|}{|y|}$。

\corepoint{近似公式的前提是 $|x|$（或 $|\Delta x|$）足够小，$\alpha$ 可为任意实数；$x$ 用弧度。用微分作近似时，$\Delta x$ 要取成“方便计算的小量”（如由 $\sqrt{4.02}$ 想到 $x_0=4$、$\Delta x=0.02$）。误差估计的常见坑：多个量相乘时相对误差要相加，不能只算一项。}
 引入，不含 \documentclass / \begin{document}

\section{导数概念}

\subsection{两个引例}
引例一（变速直线运动的速度）：位移 $s=s(t)$，从 $t_0$ 到 $t_0+\Delta t$ 的平均速度 $\bar v=\dfrac{\Delta s}{\Delta t}$，令 $\Delta t\to0$ 得瞬时速度 $v(t_0)=\lim\limits_{\Delta t\to0}\dfrac{\Delta s}{\Delta t}$。

引例二（平面曲线的切线斜率）：曲线 $y=f(x)$ 上点 $M(x_0,y_0)$ 处切线斜率 $k=\lim\limits_{\Delta x\to0}\dfrac{\Delta y}{\Delta x}$。

两个引例的形式完全相同：都是\textbf{函数增量与自变量增量之比当自变量增量趋于零时的极限}。抽去具体意义就得到导数概念。

\subsection{导数的定义}
\textbf{定义}\quad 设函数 $y=f(x)$ 在点 $x_0$ 的某个邻域内有定义，当自变量 $x$ 在 $x_0$ 处取得增量 $\Delta x$（$x_0+\Delta x$ 仍在该邻域内）时，相应地函数取得增量
\[
\Delta y=f(x_0+\Delta x)-f(x_0);
\]
如果 $\Delta y$ 与 $\Delta x$ 之比当 $\Delta x\to0$ 时的极限存在，则称函数 $y=f(x)$ 在点 $x_0$ 处\textbf{可导}，并称这个极限为函数 $y=f(x)$ 在点 $x_0$ 处的\textbf{导数}，记为 $y'|_{x=x_0}$、$\dfrac{\mathrm{d}y}{\mathrm{d}x}\Big|_{x=x_0}$ 或 $f'(x_0)$，即
\keyformula{f'(x_0)=\lim_{\Delta x\to0}\frac{\Delta y}{\Delta x}=\lim_{\Delta x\to0}\frac{f(x_0+\Delta x)-f(x_0)}{\Delta x}}
\noindent 也可写成等价形式（令 $x=x_0+\Delta x$）
\keyformula{f'(x_0)=\lim_{x\to x_0}\frac{f(x)-f(x_0)}{x-x_0}=\lim_{h\to0}\frac{f(x_0+h)-f(x_0)}{h}}
\noindent 若上述极限不存在，则称 $f(x)$ 在 $x_0$ 处不可导；若极限为无穷大，习惯上也称导数不存在（导数为无穷大）。

\textbf{导函数}\quad 若 $f(x)$ 在区间 $I$ 内每一点都可导，则对每个 $x\in I$ 都有唯一确定的导数值 $f'(x)$ 与之对应，由此得到 $I$ 上的新函数，称为 $f(x)$ 的导函数，记作 $y'$、$f'(x)$ 或 $\dfrac{\mathrm{d}y}{\mathrm{d}x}$。

\subsubsection{单侧导数}
\[\text{左导数：}f'_-(x_0)=\lim_{\Delta x\to0^-}\frac{f(x_0+\Delta x)-f(x_0)}{\Delta x};\qquad
\text{右导数：}f'_+(x_0)=\lim_{\Delta x\to0^+}\frac{f(x_0+\Delta x)-f(x_0)}{\Delta x}.\]
\keyformula{f'(x_0)\ \text{存在}\iff f'_-(x_0)=f'_+(x_0)\ \text{均存在且相等}}

\corepoint{求分段函数在\textbf{分界点}处的导数，必须用定义分左右两侧算 $f'_-$、$f'_+$ 再比较，不能用求导公式直接代——分界点两侧的表达式不同。而在各段\textbf{内部}可用公式求导。$f$ 在 $[a,b]$ 上可导指：在 $(a,b)$ 内可导，且 $f'_+(a)$、$f'_-(b)$ 存在。}

\subsection{导数的几何意义与物理意义}
$f'(x_0)$ 是曲线 $y=f(x)$ 在点 $M(x_0,y_0)$ 处\textbf{切线的斜率}。于是
\keyformula{\text{切线：}\ y-y_0=f'(x_0)(x-x_0);\qquad \text{法线：}\ y-y_0=-\frac{1}{f'(x_0)}(x-x_0)\ (f'(x_0)\neq0)}
\noindent 若 $f'(x_0)=\infty$，则切线为竖直直线 $x=x_0$。

一般地，$f'(x_0)$ 就是因变量对自变量的变化率（瞬时变化率、边际）：$v=\frac{\mathrm{d}s}{\mathrm{d}t}$ 是速度，$\rho=\frac{\mathrm{d}Q}{\mathrm{d}p}$ 是边际需求等。

\subsection{函数可导性与连续性的关系}
\textbf{定理}\quad 若函数 $y=f(x)$ 在点 $x_0$ 处可导，则它在点 $x_0$ 处必连续。

\corepoint{逆命题不成立：连续不一定可导。典型反例 $y=|x|$ 在 $x=0$ 处连续，但 $f'_-(0)=-1\neq f'_+(0)=1$，不可导。$\sqrt[3]{x}$ 在 $x=0$ 处连续且 $f'(0)=\infty$（有竖直切线）。“可导 $\Rightarrow$ 连续”常用其逆否形式：\textbf{不连续必不可导}。判断“在 $x_0$ 是否可导”的标准流程是：先看连续（不连续直接不可导）→ 再用定义算左右导数比较。}

\section{函数的求导法则}

\subsection{函数的和、差、积、商的求导法则}
设 $u=u(x)$、$v=v(x)$ 都在点 $x$ 处可导，则
\keyformula{(u\pm v)'=u'\pm v',\qquad (uv)'=u'v+uv',\qquad \Bigl(\frac uv\Bigr)'=\frac{u'v-uv'}{v^2}\ (v\neq0)}
\noindent 推论：$(Cu)'=Cu'$；$\Bigl(\dfrac{1}{v}\Bigr)'=-\dfrac{v'}{v^2}$。

\corepoint{乘积法则不能写成 $u'v'$；三个因子相乘要两项三项轮流求导：$(uvw)'=u'vw+uv'w+uvw'$。商的分子顺序是 $u'v-uv'$，与积的区别就是“减号在前”，记错会使答案差一个负号。}

\subsection{反函数的求导法则}
\textbf{定理}\quad 若函数 $x=\varphi(y)$ 在某区间内单调、可导且 $\varphi'(y)\neq0$，则它的反函数 $y=f(x)$ 在该对应区间内也可导，且
\keyformula{f'(x)=\frac{1}{\varphi'(y)}\qquad\text{即}\qquad \frac{\mathrm{d}y}{\mathrm{d}x}=\frac{1}{\frac{\mathrm{d}x}{\mathrm{d}y}}}

\corepoint{反函数的导数等于原函数导数的倒数，但两边要“换变量”：算完 $\frac{1}{\varphi'(y)}$ 后必须把 $y$ 用 $y=f(x)$ 回代成 $x$ 的表达式。这个法则主要用于\textbf{推导}反三角函数、对数函数的导数公式（因为它们的原函数好求导），而不是用来做具体计算。}

\subsection{复合函数的求导法则（链式法则）}
\textbf{定理}\quad 若 $u=g(x)$ 在点 $x$ 处可导，$y=f(u)$ 在对应点 $u=g(x)$ 处可导，则复合函数 $y=f[g(x)]$ 在点 $x$ 处可导，且
\keyformula{\frac{\mathrm{d}y}{\mathrm{d}x}=f'(u)\cdot g'(x)\qquad\text{或}\qquad y'_x=y'_u\cdot u'_x}
\noindent 多层复合逐层相乘：$y=f(u)$，$u=g(v)$，$v=h(x)$，则 $y'_x=f'(u)\,g'(v)\,h'(x)$。

\corepoint{链式法则三条铁律：(1) 不要跳层，有几层就乘几个因子；(2) 每层都“对外层变量求导，再乘内层对自变量的导数”；(3) 用 $y'_x=y'_u u'_x$ 写法时，最后要把 $u$ 回代成 $x$ 的表达式。$(\sin^2x)'=2\sin x\cos x$ 不是 $2\sin x$；$(e^{x^2})'=2xe^{x^2}$ 不要漏掉 $2x$。}

\subsection{基本求导法则与导数公式}
\begin{center}
\begin{tabular}{@{}ll@{\hspace{2em}}ll@{}}
\toprule
$f(x)$ & $f'(x)$ & $f(x)$ & $f'(x)$ \\
\midrule
$C$ & $0$ & $\sin x$ & $\cos x$ \\
$x^{\mu}$ & $\mu x^{\mu-1}$ & $\cos x$ & $-\sin x$ \\
$a^{x}$ & $a^{x}\ln a$ & $\tan x$ & $\sec^2x$ \\
$e^{x}$ & $e^{x}$ & $\cot x$ & $-\csc^2x$ \\
$\log_a x$ & $\dfrac{1}{x\ln a}$ & $\sec x$ & $\sec x\tan x$ \\
$\ln x$ & $\dfrac1x$ & $\csc x$ & $-\csc x\cot x$ \\
$\arcsin x$ & $\dfrac{1}{\sqrt{1-x^2}}$ & $\arctan x$ & $\dfrac{1}{1+x^2}$ \\
$\arccos x$ & $-\dfrac{1}{\sqrt{1-x^2}}$ & $\operatorname{arccot}x$ & $-\dfrac{1}{1+x^2}$ \\
\bottomrule
\end{tabular}
\end{center}
\noindent 熟记这张表是求一切导数的基础；配合四则法则、链式法则、反函数法则，就可求初等函数在定义区间内的任意阶导数。

\section{高阶导数}

\subsection{定义与记法}
函数 $y=f(x)$ 的导数 $f'(x)$ 仍是 $x$ 的函数。若 $f'(x)$ 在点 $x$ 处可导，则称 $f'(x)$ 的导数为 $f(x)$ 的\textbf{二阶导数}，记作
\[
y'',\quad f''(x),\quad \frac{\mathrm{d}^2y}{\mathrm{d}x^2};
\]
类似地定义三阶、四阶导数，$n$ 阶导数记作 $y^{(n)}$、$f^{(n)}(x)$、$\dfrac{\mathrm{d}^ny}{\mathrm{d}x^n}$。二阶及二阶以上的导数统称\textbf{高阶导数}。

\subsection{几个常用的 $n$ 阶导数公式}
\keyformula{\begin{aligned}
(e^{ax})^{(n)}&=a^{n}e^{ax}, &\quad (\sin x)^{(n)}&=\sin\Bigl(x+\frac{n\pi}{2}\Bigr),\\
(\cos x)^{(n)}&=\cos\Bigl(x+\frac{n\pi}{2}\Bigr), &\quad (x^{\mu})^{(n)}&=\mu(\mu-1)\cdots(\mu-n+1)x^{\mu-n},\\
(\ln x)^{(n)}&=\frac{(-1)^{n-1}(n-1)!}{x^{n}}. & &
\end{aligned}}

\corepoint{求 $n$ 阶导数的一般方法是“算前几阶、找规律、用数学归纳法证明”。$\sin x,\cos x$ 每求导一次相位前进 $\frac\pi2$，这是周期性函数的高阶导数的常用技巧；$n$ 从 $1$ 开始时 $\sin^{(n)}$ 的公式同样成立。}

\subsection{莱布尼茨公式}
若 $u=u(x)$、$v=v(x)$ 均有 $n$ 阶导数，则
\keyformula{(uv)^{(n)}=\sum_{k=0}^{n}\mathrm{C}_n^{k}\,u^{(n-k)}\,v^{(k)}}
\noindent 其中 $u^{(0)}=u$，$v^{(0)}=v$。形式与二项式展开定理完全相同。

\corepoint{莱布尼茨公式适用于“两个因子之积、其中一个因子（如 $x^m$、多项式）求若干阶导后变零”的情形：多项式 $x^m$ 的 $n$ 阶导在 $n>m$ 时为零，所以和式只剩前面有限项。用错方向的典型是把 $e^x\ln(1+x)$ 硬套莱布尼茨——那不划算。}

\section{隐函数及由参数方程所确定的函数的导数·相关变化率}

\subsection{隐函数的导数}
若方程 $F(x,y)=0$ 确定了 $y$ 是 $x$ 的函数，则称 $y$ 为 $x$ 的\textbf{隐函数}。求导方法：把 $y$ 看作 $x$ 的函数（$y=y(x)$），对方程两边同时关于 $x$ 求导，遇到 $y$ 时按链式法则乘以 $y'$，得到含 $y'$ 的方程，解出 $y'$。

\corepoint{隐函数求导的三个要点：(1) 记住 $y$ 是 $x$ 的函数，$\frac{\mathrm{d}}{\mathrm{d}x}[f(y)]=f'(y)y'$，如 $\frac{\mathrm{d}}{\mathrm{d}x}(y^2)=2yy'$；(2) 结果是含 $x,y$ 的式子，通常保留 $y$（不必解出 $y$）；(3) 求二阶导时，把一阶导的结果\textbf{代入}后再求导，不要停下。求曲线 $F(x,y)=0$ 在点 $(x_0,y_0)$ 处的切线，可用 $y'|_{(x_0,y_0)}$ 直接代入。}

\subsubsection{对数求导法}
对 $y=u(x)^{v(x)}$（$u>0$）或幂连乘、连除式，先在两边取对数 $\ln y=v\ln u$，再两边求导：
\[
\frac{y'}{y}=v'\ln u+v\frac{u'}{u}\ \Longrightarrow\ y'=u^{v}\Bigl(v'\ln u+\frac{vu'}{u}\Bigr).
\]

\corepoint{对数求导法的适用信号：幂指函数（底与指数都含 $x$）、多个因子连乘连除、含根式与乘方。它把乘除化为加减、把指数化为乘数。结果是 $y'$ 的表达式里要\textbf{回乘 $y$ 本身}，别漏掉这个因子（忘记回代是最常见错误）。}

\subsection{由参数方程所确定的函数的导数}
设 $x=\varphi(t)$、$y=\psi(t)$ 都可导且 $\varphi'(t)\neq0$，则
\keyformula{\frac{\mathrm{d}y}{\mathrm{d}x}=\frac{\frac{\mathrm{d}y}{\mathrm{d}t}}{\frac{\mathrm{d}x}{\mathrm{d}t}}=\frac{\psi'(t)}{\varphi'(t)}}
\noindent 二阶导数（对 $x$ 再求一次导，注意 $x$ 仍是 $t$ 的函数）：
\keyformula{\frac{\mathrm{d}^2y}{\mathrm{d}x^2}=\frac{\mathrm{d}}{\mathrm{d}t}\Bigl(\frac{\mathrm{d}y}{\mathrm{d}x}\Bigr)\cdot\frac{1}{\varphi'(t)}
=\frac{\psi''(t)\varphi'(t)-\psi'(t)\varphi''(t)}{\bigl[\varphi'(t)\bigr]^{3}}}

\corepoint{二阶导数\textbf{绝不能}写成 $\dfrac{\psi''(t)}{\varphi''(t)}$。正确做法：先把一阶导 $\dfrac{\mathrm{d}y}{\mathrm{d}x}$ 看成 $t$ 的函数，对它关于 $t$ 求导，再乘以 $\dfrac{\mathrm{d}t}{\mathrm{d}x}=\dfrac{1}{\varphi'(t)}$。这一条几乎每份卷子都在考。}

\subsection{相关变化率}
设 $x=x(t)$、$y=y(t)$ 都可导，而 $x$ 与 $y$ 之间存在关系 $F(x,y)=0$。对 $t$ 求导得
\[
\frac{\mathrm{d}F}{\mathrm{d}x}\cdot\frac{\mathrm{d}x}{\mathrm{d}t}+\frac{\mathrm{d}F}{\mathrm{d}y}\cdot\frac{\mathrm{d}y}{\mathrm{d}t}=0,
\]
从而可由一个变化率求出另一个变化率。这类问题称为\textbf{相关变化率}问题。

\corepoint{解相关变化率问题的步骤：(1) 建立变量之间的方程（几何关系、物理定律）；(2) 两边同时对 $t$ 求导（所有变量都是 $t$ 的函数）；(3) 代入某一时刻的已知值求出待求变化率。注意“代入数值”必须放在求导\textbf{之后}——先代入会把变量变成常数而使导数为 0。}

\section{函数的微分}

\subsection{微分的定义}
\textbf{定义}\quad 设函数 $y=f(x)$ 在某区间内有定义，$x_0$ 及 $x_0+\Delta x$ 都在该区间内。如果函数的增量可表示为
\[
\Delta y=f(x_0+\Delta x)-f(x_0)=A\,\Delta x+o(\Delta x),
\]
其中 $A$ 是不依赖于 $\Delta x$ 的常数，则称 $f(x)$ 在点 $x_0$ 处\textbf{可微}，而 $A\,\Delta x$ 称为 $f(x)$ 在点 $x_0$ 处的\textbf{微分}，记作 $\mathrm{d}y$，即 $\mathrm{d}y=A\,\Delta x$。

\keyformula{f(x)\ \text{在}\ x_0\ \text{处可微}\iff f(x)\ \text{在}\ x_0\ \text{处可导，且}\ \mathrm{d}y=f'(x_0)\,\Delta x}

\noindent 又因 $\mathrm{d}x=\Delta x$（自变量的微分等于自变量的增量），所以
\keyformula{\mathrm{d}y=f'(x)\,\mathrm{d}x,\qquad \frac{\mathrm{d}y}{\mathrm{d}x}=f'(x)}
\noindent 后一式说明：\textbf{导数就是函数的微分与自变量的微分之商}，这就是“微商”一词的来历。

\corepoint{一元函数“可微 $\iff$ 可导”，二者是同一件事；但 $\Delta y$ 与 $\mathrm{d}y$ 不同：$\Delta y$ 是真正的函数增量，$\mathrm{d}y=f'(x)\Delta x$ 是它的\textbf{线性主部}，$\Delta y=\mathrm{d}y+o(\Delta x)$。$\mathrm{d}y$ 依赖于 $x$ 与 $\Delta x$ 两个量，$\mathrm{d}x$ 与 $\Delta x$ 恒等。求 $\mathrm{d}y$ 时不要忘记最后的 $\mathrm{d}x$。}

\subsection{微分的几何意义}
在曲线 $y=f(x)$ 上取点 $M(x_0,y_0)$，当横坐标增加 $\Delta x$ 时，曲线上的纵坐标增量为 $\Delta y=NT$，而切线上相应的纵坐标增量为 $\mathrm{d}y$。因此 $\mathrm{d}y$ 是曲线在点 $M$ 处\textbf{切线的纵坐标增量}，即“以直代曲”的线性近似。

\subsection{基本初等函数的微分公式与微分运算法则}
把求导公式两边同乘 $\mathrm{d}x$ 即得微分公式，如
\[
\mathrm{d}(x^{\mu})=\mu x^{\mu-1}\mathrm{d}x,\quad \mathrm{d}(\sin x)=\cos x\,\mathrm{d}x,\quad
\mathrm{d}(\arctan x)=\frac{\mathrm{d}x}{1+x^2},\quad \mathrm{d}(\ln x)=\frac{\mathrm{d}x}{x}.
\]
微分运算法则：
\keyformula{\mathrm{d}(u\pm v)=\mathrm{d}u\pm\mathrm{d}v,\qquad \mathrm{d}(uv)=v\,\mathrm{d}u+u\,\mathrm{d}v,\qquad \mathrm{d}\Bigl(\frac uv\Bigr)=\frac{v\,\mathrm{d}u-u\,\mathrm{d}v}{v^{2}}}

\subsubsection{复合函数的微分与一阶微分形式不变性}
设 $y=f(u)$ 可微，$u=g(x)$ 可微，则 $\mathrm{d}y=f'(u)\,\mathrm{d}u$。注意这里 $u$ 既可以是自变量，也可以是另一个变量的函数，形式都是 $f'(u)\mathrm{d}u$，这就是\textbf{一阶微分形式不变性}。

\corepoint{一阶微分形式不变性指的是“形式”不变（$\mathrm{d}y=f'(u)\mathrm{d}u$ 恒成立），不是说导数不变（$y'_x=y'_u u'_x$ 差一个 $u'_x$）。它使凑微分（如 $\cos2x\,\mathrm{d}x=\frac12\mathrm{d}(\sin2x)$）成为可能，是第 4 章换元积分法的基础。}

\subsection{微分在近似计算中的应用}
当 $|\Delta x|$ 很小时，
\keyformula{\Delta y\approx \mathrm{d}y=f'(x_0)\Delta x,\qquad f(x_0+\Delta x)\approx f(x_0)+f'(x_0)\Delta x}
\noindent 特别地，取 $x_0=0$ 得 $f(x)\approx f(0)+f'(0)x$，由此得到常用近似式：
\keyformula{(1+x)^{\alpha}\approx1+\alpha x,\quad \sin x\approx x,\quad \tan x\approx x,\quad e^{x}\approx1+x,\quad \ln(1+x)\approx x\quad(|x|\ \text{很小})}

\noindent 误差估计：绝对误差 $|\Delta y|\approx|\mathrm{d}y|=|f'(x)|\,\Delta x$，相对误差 $\dfrac{|\mathrm{d}y|}{|y|}$。

\corepoint{近似公式的前提是 $|x|$（或 $|\Delta x|$）足够小，$\alpha$ 可为任意实数；$x$ 用弧度。用微分作近似时，$\Delta x$ 要取成“方便计算的小量”（如由 $\sqrt{4.02}$ 想到 $x_0=4$、$\Delta x=0.02$）。误差估计的常见坑：多个量相乘时相对误差要相加，不能只算一项。}
 引入，不含 \documentclass / \begin{document}

\section{导数概念}

\subsection{两个引例}
引例一（变速直线运动的速度）：位移 $s=s(t)$，从 $t_0$ 到 $t_0+\Delta t$ 的平均速度 $\bar v=\dfrac{\Delta s}{\Delta t}$，令 $\Delta t\to0$ 得瞬时速度 $v(t_0)=\lim\limits_{\Delta t\to0}\dfrac{\Delta s}{\Delta t}$。

引例二（平面曲线的切线斜率）：曲线 $y=f(x)$ 上点 $M(x_0,y_0)$ 处切线斜率 $k=\lim\limits_{\Delta x\to0}\dfrac{\Delta y}{\Delta x}$。

两个引例的形式完全相同：都是\textbf{函数增量与自变量增量之比当自变量增量趋于零时的极限}。抽去具体意义就得到导数概念。

\subsection{导数的定义}
\textbf{定义}\quad 设函数 $y=f(x)$ 在点 $x_0$ 的某个邻域内有定义，当自变量 $x$ 在 $x_0$ 处取得增量 $\Delta x$（$x_0+\Delta x$ 仍在该邻域内）时，相应地函数取得增量
\[
\Delta y=f(x_0+\Delta x)-f(x_0);
\]
如果 $\Delta y$ 与 $\Delta x$ 之比当 $\Delta x\to0$ 时的极限存在，则称函数 $y=f(x)$ 在点 $x_0$ 处\textbf{可导}，并称这个极限为函数 $y=f(x)$ 在点 $x_0$ 处的\textbf{导数}，记为 $y'|_{x=x_0}$、$\dfrac{\mathrm{d}y}{\mathrm{d}x}\Big|_{x=x_0}$ 或 $f'(x_0)$，即
\keyformula{f'(x_0)=\lim_{\Delta x\to0}\frac{\Delta y}{\Delta x}=\lim_{\Delta x\to0}\frac{f(x_0+\Delta x)-f(x_0)}{\Delta x}}
\noindent 也可写成等价形式（令 $x=x_0+\Delta x$）
\keyformula{f'(x_0)=\lim_{x\to x_0}\frac{f(x)-f(x_0)}{x-x_0}=\lim_{h\to0}\frac{f(x_0+h)-f(x_0)}{h}}
\noindent 若上述极限不存在，则称 $f(x)$ 在 $x_0$ 处不可导；若极限为无穷大，习惯上也称导数不存在（导数为无穷大）。

\textbf{导函数}\quad 若 $f(x)$ 在区间 $I$ 内每一点都可导，则对每个 $x\in I$ 都有唯一确定的导数值 $f'(x)$ 与之对应，由此得到 $I$ 上的新函数，称为 $f(x)$ 的导函数，记作 $y'$、$f'(x)$ 或 $\dfrac{\mathrm{d}y}{\mathrm{d}x}$。

\subsubsection{单侧导数}
\[\text{左导数：}f'_-(x_0)=\lim_{\Delta x\to0^-}\frac{f(x_0+\Delta x)-f(x_0)}{\Delta x};\qquad
\text{右导数：}f'_+(x_0)=\lim_{\Delta x\to0^+}\frac{f(x_0+\Delta x)-f(x_0)}{\Delta x}.\]
\keyformula{f'(x_0)\ \text{存在}\iff f'_-(x_0)=f'_+(x_0)\ \text{均存在且相等}}

\corepoint{求分段函数在\textbf{分界点}处的导数，必须用定义分左右两侧算 $f'_-$、$f'_+$ 再比较，不能用求导公式直接代——分界点两侧的表达式不同。而在各段\textbf{内部}可用公式求导。$f$ 在 $[a,b]$ 上可导指：在 $(a,b)$ 内可导，且 $f'_+(a)$、$f'_-(b)$ 存在。}

\subsection{导数的几何意义与物理意义}
$f'(x_0)$ 是曲线 $y=f(x)$ 在点 $M(x_0,y_0)$ 处\textbf{切线的斜率}。于是
\keyformula{\text{切线：}\ y-y_0=f'(x_0)(x-x_0);\qquad \text{法线：}\ y-y_0=-\frac{1}{f'(x_0)}(x-x_0)\ (f'(x_0)\neq0)}
\noindent 若 $f'(x_0)=\infty$，则切线为竖直直线 $x=x_0$。

一般地，$f'(x_0)$ 就是因变量对自变量的变化率（瞬时变化率、边际）：$v=\frac{\mathrm{d}s}{\mathrm{d}t}$ 是速度，$\rho=\frac{\mathrm{d}Q}{\mathrm{d}p}$ 是边际需求等。

\subsection{函数可导性与连续性的关系}
\textbf{定理}\quad 若函数 $y=f(x)$ 在点 $x_0$ 处可导，则它在点 $x_0$ 处必连续。

\corepoint{逆命题不成立：连续不一定可导。典型反例 $y=|x|$ 在 $x=0$ 处连续，但 $f'_-(0)=-1\neq f'_+(0)=1$，不可导。$\sqrt[3]{x}$ 在 $x=0$ 处连续且 $f'(0)=\infty$（有竖直切线）。“可导 $\Rightarrow$ 连续”常用其逆否形式：\textbf{不连续必不可导}。判断“在 $x_0$ 是否可导”的标准流程是：先看连续（不连续直接不可导）→ 再用定义算左右导数比较。}

\section{函数的求导法则}

\subsection{函数的和、差、积、商的求导法则}
设 $u=u(x)$、$v=v(x)$ 都在点 $x$ 处可导，则
\keyformula{(u\pm v)'=u'\pm v',\qquad (uv)'=u'v+uv',\qquad \Bigl(\frac uv\Bigr)'=\frac{u'v-uv'}{v^2}\ (v\neq0)}
\noindent 推论：$(Cu)'=Cu'$；$\Bigl(\dfrac{1}{v}\Bigr)'=-\dfrac{v'}{v^2}$。

\corepoint{乘积法则不能写成 $u'v'$；三个因子相乘要两项三项轮流求导：$(uvw)'=u'vw+uv'w+uvw'$。商的分子顺序是 $u'v-uv'$，与积的区别就是“减号在前”，记错会使答案差一个负号。}

\subsection{反函数的求导法则}
\textbf{定理}\quad 若函数 $x=\varphi(y)$ 在某区间内单调、可导且 $\varphi'(y)\neq0$，则它的反函数 $y=f(x)$ 在该对应区间内也可导，且
\keyformula{f'(x)=\frac{1}{\varphi'(y)}\qquad\text{即}\qquad \frac{\mathrm{d}y}{\mathrm{d}x}=\frac{1}{\frac{\mathrm{d}x}{\mathrm{d}y}}}

\corepoint{反函数的导数等于原函数导数的倒数，但两边要“换变量”：算完 $\frac{1}{\varphi'(y)}$ 后必须把 $y$ 用 $y=f(x)$ 回代成 $x$ 的表达式。这个法则主要用于\textbf{推导}反三角函数、对数函数的导数公式（因为它们的原函数好求导），而不是用来做具体计算。}

\subsection{复合函数的求导法则（链式法则）}
\textbf{定理}\quad 若 $u=g(x)$ 在点 $x$ 处可导，$y=f(u)$ 在对应点 $u=g(x)$ 处可导，则复合函数 $y=f[g(x)]$ 在点 $x$ 处可导，且
\keyformula{\frac{\mathrm{d}y}{\mathrm{d}x}=f'(u)\cdot g'(x)\qquad\text{或}\qquad y'_x=y'_u\cdot u'_x}
\noindent 多层复合逐层相乘：$y=f(u)$，$u=g(v)$，$v=h(x)$，则 $y'_x=f'(u)\,g'(v)\,h'(x)$。

\corepoint{链式法则三条铁律：(1) 不要跳层，有几层就乘几个因子；(2) 每层都“对外层变量求导，再乘内层对自变量的导数”；(3) 用 $y'_x=y'_u u'_x$ 写法时，最后要把 $u$ 回代成 $x$ 的表达式。$(\sin^2x)'=2\sin x\cos x$ 不是 $2\sin x$；$(e^{x^2})'=2xe^{x^2}$ 不要漏掉 $2x$。}

\subsection{基本求导法则与导数公式}
\begin{center}
\begin{tabular}{@{}ll@{\hspace{2em}}ll@{}}
\toprule
$f(x)$ & $f'(x)$ & $f(x)$ & $f'(x)$ \\
\midrule
$C$ & $0$ & $\sin x$ & $\cos x$ \\
$x^{\mu}$ & $\mu x^{\mu-1}$ & $\cos x$ & $-\sin x$ \\
$a^{x}$ & $a^{x}\ln a$ & $\tan x$ & $\sec^2x$ \\
$e^{x}$ & $e^{x}$ & $\cot x$ & $-\csc^2x$ \\
$\log_a x$ & $\dfrac{1}{x\ln a}$ & $\sec x$ & $\sec x\tan x$ \\
$\ln x$ & $\dfrac1x$ & $\csc x$ & $-\csc x\cot x$ \\
$\arcsin x$ & $\dfrac{1}{\sqrt{1-x^2}}$ & $\arctan x$ & $\dfrac{1}{1+x^2}$ \\
$\arccos x$ & $-\dfrac{1}{\sqrt{1-x^2}}$ & $\operatorname{arccot}x$ & $-\dfrac{1}{1+x^2}$ \\
\bottomrule
\end{tabular}
\end{center}
\noindent 熟记这张表是求一切导数的基础；配合四则法则、链式法则、反函数法则，就可求初等函数在定义区间内的任意阶导数。

\section{高阶导数}

\subsection{定义与记法}
函数 $y=f(x)$ 的导数 $f'(x)$ 仍是 $x$ 的函数。若 $f'(x)$ 在点 $x$ 处可导，则称 $f'(x)$ 的导数为 $f(x)$ 的\textbf{二阶导数}，记作
\[
y'',\quad f''(x),\quad \frac{\mathrm{d}^2y}{\mathrm{d}x^2};
\]
类似地定义三阶、四阶导数，$n$ 阶导数记作 $y^{(n)}$、$f^{(n)}(x)$、$\dfrac{\mathrm{d}^ny}{\mathrm{d}x^n}$。二阶及二阶以上的导数统称\textbf{高阶导数}。

\subsection{几个常用的 $n$ 阶导数公式}
\keyformula{\begin{aligned}
(e^{ax})^{(n)}&=a^{n}e^{ax}, &\quad (\sin x)^{(n)}&=\sin\Bigl(x+\frac{n\pi}{2}\Bigr),\\
(\cos x)^{(n)}&=\cos\Bigl(x+\frac{n\pi}{2}\Bigr), &\quad (x^{\mu})^{(n)}&=\mu(\mu-1)\cdots(\mu-n+1)x^{\mu-n},\\
(\ln x)^{(n)}&=\frac{(-1)^{n-1}(n-1)!}{x^{n}}. & &
\end{aligned}}

\corepoint{求 $n$ 阶导数的一般方法是“算前几阶、找规律、用数学归纳法证明”。$\sin x,\cos x$ 每求导一次相位前进 $\frac\pi2$，这是周期性函数的高阶导数的常用技巧；$n$ 从 $1$ 开始时 $\sin^{(n)}$ 的公式同样成立。}

\subsection{莱布尼茨公式}
若 $u=u(x)$、$v=v(x)$ 均有 $n$ 阶导数，则
\keyformula{(uv)^{(n)}=\sum_{k=0}^{n}\mathrm{C}_n^{k}\,u^{(n-k)}\,v^{(k)}}
\noindent 其中 $u^{(0)}=u$，$v^{(0)}=v$。形式与二项式展开定理完全相同。

\corepoint{莱布尼茨公式适用于“两个因子之积、其中一个因子（如 $x^m$、多项式）求若干阶导后变零”的情形：多项式 $x^m$ 的 $n$ 阶导在 $n>m$ 时为零，所以和式只剩前面有限项。用错方向的典型是把 $e^x\ln(1+x)$ 硬套莱布尼茨——那不划算。}

\section{隐函数及由参数方程所确定的函数的导数·相关变化率}

\subsection{隐函数的导数}
若方程 $F(x,y)=0$ 确定了 $y$ 是 $x$ 的函数，则称 $y$ 为 $x$ 的\textbf{隐函数}。求导方法：把 $y$ 看作 $x$ 的函数（$y=y(x)$），对方程两边同时关于 $x$ 求导，遇到 $y$ 时按链式法则乘以 $y'$，得到含 $y'$ 的方程，解出 $y'$。

\corepoint{隐函数求导的三个要点：(1) 记住 $y$ 是 $x$ 的函数，$\frac{\mathrm{d}}{\mathrm{d}x}[f(y)]=f'(y)y'$，如 $\frac{\mathrm{d}}{\mathrm{d}x}(y^2)=2yy'$；(2) 结果是含 $x,y$ 的式子，通常保留 $y$（不必解出 $y$）；(3) 求二阶导时，把一阶导的结果\textbf{代入}后再求导，不要停下。求曲线 $F(x,y)=0$ 在点 $(x_0,y_0)$ 处的切线，可用 $y'|_{(x_0,y_0)}$ 直接代入。}

\subsubsection{对数求导法}
对 $y=u(x)^{v(x)}$（$u>0$）或幂连乘、连除式，先在两边取对数 $\ln y=v\ln u$，再两边求导：
\[
\frac{y'}{y}=v'\ln u+v\frac{u'}{u}\ \Longrightarrow\ y'=u^{v}\Bigl(v'\ln u+\frac{vu'}{u}\Bigr).
\]

\corepoint{对数求导法的适用信号：幂指函数（底与指数都含 $x$）、多个因子连乘连除、含根式与乘方。它把乘除化为加减、把指数化为乘数。结果是 $y'$ 的表达式里要\textbf{回乘 $y$ 本身}，别漏掉这个因子（忘记回代是最常见错误）。}

\subsection{由参数方程所确定的函数的导数}
设 $x=\varphi(t)$、$y=\psi(t)$ 都可导且 $\varphi'(t)\neq0$，则
\keyformula{\frac{\mathrm{d}y}{\mathrm{d}x}=\frac{\frac{\mathrm{d}y}{\mathrm{d}t}}{\frac{\mathrm{d}x}{\mathrm{d}t}}=\frac{\psi'(t)}{\varphi'(t)}}
\noindent 二阶导数（对 $x$ 再求一次导，注意 $x$ 仍是 $t$ 的函数）：
\keyformula{\frac{\mathrm{d}^2y}{\mathrm{d}x^2}=\frac{\mathrm{d}}{\mathrm{d}t}\Bigl(\frac{\mathrm{d}y}{\mathrm{d}x}\Bigr)\cdot\frac{1}{\varphi'(t)}
=\frac{\psi''(t)\varphi'(t)-\psi'(t)\varphi''(t)}{\bigl[\varphi'(t)\bigr]^{3}}}

\corepoint{二阶导数\textbf{绝不能}写成 $\dfrac{\psi''(t)}{\varphi''(t)}$。正确做法：先把一阶导 $\dfrac{\mathrm{d}y}{\mathrm{d}x}$ 看成 $t$ 的函数，对它关于 $t$ 求导，再乘以 $\dfrac{\mathrm{d}t}{\mathrm{d}x}=\dfrac{1}{\varphi'(t)}$。这一条几乎每份卷子都在考。}

\subsection{相关变化率}
设 $x=x(t)$、$y=y(t)$ 都可导，而 $x$ 与 $y$ 之间存在关系 $F(x,y)=0$。对 $t$ 求导得
\[
\frac{\mathrm{d}F}{\mathrm{d}x}\cdot\frac{\mathrm{d}x}{\mathrm{d}t}+\frac{\mathrm{d}F}{\mathrm{d}y}\cdot\frac{\mathrm{d}y}{\mathrm{d}t}=0,
\]
从而可由一个变化率求出另一个变化率。这类问题称为\textbf{相关变化率}问题。

\corepoint{解相关变化率问题的步骤：(1) 建立变量之间的方程（几何关系、物理定律）；(2) 两边同时对 $t$ 求导（所有变量都是 $t$ 的函数）；(3) 代入某一时刻的已知值求出待求变化率。注意“代入数值”必须放在求导\textbf{之后}——先代入会把变量变成常数而使导数为 0。}

\section{函数的微分}

\subsection{微分的定义}
\textbf{定义}\quad 设函数 $y=f(x)$ 在某区间内有定义，$x_0$ 及 $x_0+\Delta x$ 都在该区间内。如果函数的增量可表示为
\[
\Delta y=f(x_0+\Delta x)-f(x_0)=A\,\Delta x+o(\Delta x),
\]
其中 $A$ 是不依赖于 $\Delta x$ 的常数，则称 $f(x)$ 在点 $x_0$ 处\textbf{可微}，而 $A\,\Delta x$ 称为 $f(x)$ 在点 $x_0$ 处的\textbf{微分}，记作 $\mathrm{d}y$，即 $\mathrm{d}y=A\,\Delta x$。

\keyformula{f(x)\ \text{在}\ x_0\ \text{处可微}\iff f(x)\ \text{在}\ x_0\ \text{处可导，且}\ \mathrm{d}y=f'(x_0)\,\Delta x}

\noindent 又因 $\mathrm{d}x=\Delta x$（自变量的微分等于自变量的增量），所以
\keyformula{\mathrm{d}y=f'(x)\,\mathrm{d}x,\qquad \frac{\mathrm{d}y}{\mathrm{d}x}=f'(x)}
\noindent 后一式说明：\textbf{导数就是函数的微分与自变量的微分之商}，这就是“微商”一词的来历。

\corepoint{一元函数“可微 $\iff$ 可导”，二者是同一件事；但 $\Delta y$ 与 $\mathrm{d}y$ 不同：$\Delta y$ 是真正的函数增量，$\mathrm{d}y=f'(x)\Delta x$ 是它的\textbf{线性主部}，$\Delta y=\mathrm{d}y+o(\Delta x)$。$\mathrm{d}y$ 依赖于 $x$ 与 $\Delta x$ 两个量，$\mathrm{d}x$ 与 $\Delta x$ 恒等。求 $\mathrm{d}y$ 时不要忘记最后的 $\mathrm{d}x$。}

\subsection{微分的几何意义}
在曲线 $y=f(x)$ 上取点 $M(x_0,y_0)$，当横坐标增加 $\Delta x$ 时，曲线上的纵坐标增量为 $\Delta y=NT$，而切线上相应的纵坐标增量为 $\mathrm{d}y$。因此 $\mathrm{d}y$ 是曲线在点 $M$ 处\textbf{切线的纵坐标增量}，即“以直代曲”的线性近似。

\subsection{基本初等函数的微分公式与微分运算法则}
把求导公式两边同乘 $\mathrm{d}x$ 即得微分公式，如
\[
\mathrm{d}(x^{\mu})=\mu x^{\mu-1}\mathrm{d}x,\quad \mathrm{d}(\sin x)=\cos x\,\mathrm{d}x,\quad
\mathrm{d}(\arctan x)=\frac{\mathrm{d}x}{1+x^2},\quad \mathrm{d}(\ln x)=\frac{\mathrm{d}x}{x}.
\]
微分运算法则：
\keyformula{\mathrm{d}(u\pm v)=\mathrm{d}u\pm\mathrm{d}v,\qquad \mathrm{d}(uv)=v\,\mathrm{d}u+u\,\mathrm{d}v,\qquad \mathrm{d}\Bigl(\frac uv\Bigr)=\frac{v\,\mathrm{d}u-u\,\mathrm{d}v}{v^{2}}}

\subsubsection{复合函数的微分与一阶微分形式不变性}
设 $y=f(u)$ 可微，$u=g(x)$ 可微，则 $\mathrm{d}y=f'(u)\,\mathrm{d}u$。注意这里 $u$ 既可以是自变量，也可以是另一个变量的函数，形式都是 $f'(u)\mathrm{d}u$，这就是\textbf{一阶微分形式不变性}。

\corepoint{一阶微分形式不变性指的是“形式”不变（$\mathrm{d}y=f'(u)\mathrm{d}u$ 恒成立），不是说导数不变（$y'_x=y'_u u'_x$ 差一个 $u'_x$）。它使凑微分（如 $\cos2x\,\mathrm{d}x=\frac12\mathrm{d}(\sin2x)$）成为可能，是第 4 章换元积分法的基础。}

\subsection{微分在近似计算中的应用}
当 $|\Delta x|$ 很小时，
\keyformula{\Delta y\approx \mathrm{d}y=f'(x_0)\Delta x,\qquad f(x_0+\Delta x)\approx f(x_0)+f'(x_0)\Delta x}
\noindent 特别地，取 $x_0=0$ 得 $f(x)\approx f(0)+f'(0)x$，由此得到常用近似式：
\keyformula{(1+x)^{\alpha}\approx1+\alpha x,\quad \sin x\approx x,\quad \tan x\approx x,\quad e^{x}\approx1+x,\quad \ln(1+x)\approx x\quad(|x|\ \text{很小})}

\noindent 误差估计：绝对误差 $|\Delta y|\approx|\mathrm{d}y|=|f'(x)|\,\Delta x$，相对误差 $\dfrac{|\mathrm{d}y|}{|y|}$。

\corepoint{近似公式的前提是 $|x|$（或 $|\Delta x|$）足够小，$\alpha$ 可为任意实数；$x$ 用弧度。用微分作近似时，$\Delta x$ 要取成“方便计算的小量”（如由 $\sqrt{4.02}$ 想到 $x_0=4$、$\Delta x=0.02$）。误差估计的常见坑：多个量相乘时相对误差要相加，不能只算一项。}
 引入，不含 \documentclass / \begin{document}

\section{导数概念}

\subsection{两个引例}
引例一（变速直线运动的速度）：位移 $s=s(t)$，从 $t_0$ 到 $t_0+\Delta t$ 的平均速度 $\bar v=\dfrac{\Delta s}{\Delta t}$，令 $\Delta t\to0$ 得瞬时速度 $v(t_0)=\lim\limits_{\Delta t\to0}\dfrac{\Delta s}{\Delta t}$。

引例二（平面曲线的切线斜率）：曲线 $y=f(x)$ 上点 $M(x_0,y_0)$ 处切线斜率 $k=\lim\limits_{\Delta x\to0}\dfrac{\Delta y}{\Delta x}$。

两个引例的形式完全相同：都是\textbf{函数增量与自变量增量之比当自变量增量趋于零时的极限}。抽去具体意义就得到导数概念。

\subsection{导数的定义}
\textbf{定义}\quad 设函数 $y=f(x)$ 在点 $x_0$ 的某个邻域内有定义，当自变量 $x$ 在 $x_0$ 处取得增量 $\Delta x$（$x_0+\Delta x$ 仍在该邻域内）时，相应地函数取得增量
\[
\Delta y=f(x_0+\Delta x)-f(x_0);
\]
如果 $\Delta y$ 与 $\Delta x$ 之比当 $\Delta x\to0$ 时的极限存在，则称函数 $y=f(x)$ 在点 $x_0$ 处\textbf{可导}，并称这个极限为函数 $y=f(x)$ 在点 $x_0$ 处的\textbf{导数}，记为 $y'|_{x=x_0}$、$\dfrac{\mathrm{d}y}{\mathrm{d}x}\Big|_{x=x_0}$ 或 $f'(x_0)$，即
\keyformula{f'(x_0)=\lim_{\Delta x\to0}\frac{\Delta y}{\Delta x}=\lim_{\Delta x\to0}\frac{f(x_0+\Delta x)-f(x_0)}{\Delta x}}
\noindent 也可写成等价形式（令 $x=x_0+\Delta x$）
\keyformula{f'(x_0)=\lim_{x\to x_0}\frac{f(x)-f(x_0)}{x-x_0}=\lim_{h\to0}\frac{f(x_0+h)-f(x_0)}{h}}
\noindent 若上述极限不存在，则称 $f(x)$ 在 $x_0$ 处不可导；若极限为无穷大，习惯上也称导数不存在（导数为无穷大）。

\textbf{导函数}\quad 若 $f(x)$ 在区间 $I$ 内每一点都可导，则对每个 $x\in I$ 都有唯一确定的导数值 $f'(x)$ 与之对应，由此得到 $I$ 上的新函数，称为 $f(x)$ 的导函数，记作 $y'$、$f'(x)$ 或 $\dfrac{\mathrm{d}y}{\mathrm{d}x}$。

\subsubsection{单侧导数}
\[\text{左导数：}f'_-(x_0)=\lim_{\Delta x\to0^-}\frac{f(x_0+\Delta x)-f(x_0)}{\Delta x};\qquad
\text{右导数：}f'_+(x_0)=\lim_{\Delta x\to0^+}\frac{f(x_0+\Delta x)-f(x_0)}{\Delta x}.\]
\keyformula{f'(x_0)\ \text{存在}\iff f'_-(x_0)=f'_+(x_0)\ \text{均存在且相等}}

\corepoint{求分段函数在\textbf{分界点}处的导数，必须用定义分左右两侧算 $f'_-$、$f'_+$ 再比较，不能用求导公式直接代——分界点两侧的表达式不同。而在各段\textbf{内部}可用公式求导。$f$ 在 $[a,b]$ 上可导指：在 $(a,b)$ 内可导，且 $f'_+(a)$、$f'_-(b)$ 存在。}

\subsection{导数的几何意义与物理意义}
$f'(x_0)$ 是曲线 $y=f(x)$ 在点 $M(x_0,y_0)$ 处\textbf{切线的斜率}。于是
\keyformula{\text{切线：}\ y-y_0=f'(x_0)(x-x_0);\qquad \text{法线：}\ y-y_0=-\frac{1}{f'(x_0)}(x-x_0)\ (f'(x_0)\neq0)}
\noindent 若 $f'(x_0)=\infty$，则切线为竖直直线 $x=x_0$。

一般地，$f'(x_0)$ 就是因变量对自变量的变化率（瞬时变化率、边际）：$v=\frac{\mathrm{d}s}{\mathrm{d}t}$ 是速度，$\rho=\frac{\mathrm{d}Q}{\mathrm{d}p}$ 是边际需求等。

\subsection{函数可导性与连续性的关系}
\textbf{定理}\quad 若函数 $y=f(x)$ 在点 $x_0$ 处可导，则它在点 $x_0$ 处必连续。

\corepoint{逆命题不成立：连续不一定可导。典型反例 $y=|x|$ 在 $x=0$ 处连续，但 $f'_-(0)=-1\neq f'_+(0)=1$，不可导。$\sqrt[3]{x}$ 在 $x=0$ 处连续且 $f'(0)=\infty$（有竖直切线）。“可导 $\Rightarrow$ 连续”常用其逆否形式：\textbf{不连续必不可导}。判断“在 $x_0$ 是否可导”的标准流程是：先看连续（不连续直接不可导）→ 再用定义算左右导数比较。}

\section{函数的求导法则}

\subsection{函数的和、差、积、商的求导法则}
设 $u=u(x)$、$v=v(x)$ 都在点 $x$ 处可导，则
\keyformula{(u\pm v)'=u'\pm v',\qquad (uv)'=u'v+uv',\qquad \Bigl(\frac uv\Bigr)'=\frac{u'v-uv'}{v^2}\ (v\neq0)}
\noindent 推论：$(Cu)'=Cu'$；$\Bigl(\dfrac{1}{v}\Bigr)'=-\dfrac{v'}{v^2}$。

\corepoint{乘积法则不能写成 $u'v'$；三个因子相乘要两项三项轮流求导：$(uvw)'=u'vw+uv'w+uvw'$。商的分子顺序是 $u'v-uv'$，与积的区别就是“减号在前”，记错会使答案差一个负号。}

\subsection{反函数的求导法则}
\textbf{定理}\quad 若函数 $x=\varphi(y)$ 在某区间内单调、可导且 $\varphi'(y)\neq0$，则它的反函数 $y=f(x)$ 在该对应区间内也可导，且
\keyformula{f'(x)=\frac{1}{\varphi'(y)}\qquad\text{即}\qquad \frac{\mathrm{d}y}{\mathrm{d}x}=\frac{1}{\frac{\mathrm{d}x}{\mathrm{d}y}}}

\corepoint{反函数的导数等于原函数导数的倒数，但两边要“换变量”：算完 $\frac{1}{\varphi'(y)}$ 后必须把 $y$ 用 $y=f(x)$ 回代成 $x$ 的表达式。这个法则主要用于\textbf{推导}反三角函数、对数函数的导数公式（因为它们的原函数好求导），而不是用来做具体计算。}

\subsection{复合函数的求导法则（链式法则）}
\textbf{定理}\quad 若 $u=g(x)$ 在点 $x$ 处可导，$y=f(u)$ 在对应点 $u=g(x)$ 处可导，则复合函数 $y=f[g(x)]$ 在点 $x$ 处可导，且
\keyformula{\frac{\mathrm{d}y}{\mathrm{d}x}=f'(u)\cdot g'(x)\qquad\text{或}\qquad y'_x=y'_u\cdot u'_x}
\noindent 多层复合逐层相乘：$y=f(u)$，$u=g(v)$，$v=h(x)$，则 $y'_x=f'(u)\,g'(v)\,h'(x)$。

\corepoint{链式法则三条铁律：(1) 不要跳层，有几层就乘几个因子；(2) 每层都“对外层变量求导，再乘内层对自变量的导数”；(3) 用 $y'_x=y'_u u'_x$ 写法时，最后要把 $u$ 回代成 $x$ 的表达式。$(\sin^2x)'=2\sin x\cos x$ 不是 $2\sin x$；$(e^{x^2})'=2xe^{x^2}$ 不要漏掉 $2x$。}

\subsection{基本求导法则与导数公式}
\begin{center}
\begin{tabular}{@{}ll@{\hspace{2em}}ll@{}}
\toprule
$f(x)$ & $f'(x)$ & $f(x)$ & $f'(x)$ \\
\midrule
$C$ & $0$ & $\sin x$ & $\cos x$ \\
$x^{\mu}$ & $\mu x^{\mu-1}$ & $\cos x$ & $-\sin x$ \\
$a^{x}$ & $a^{x}\ln a$ & $\tan x$ & $\sec^2x$ \\
$e^{x}$ & $e^{x}$ & $\cot x$ & $-\csc^2x$ \\
$\log_a x$ & $\dfrac{1}{x\ln a}$ & $\sec x$ & $\sec x\tan x$ \\
$\ln x$ & $\dfrac1x$ & $\csc x$ & $-\csc x\cot x$ \\
$\arcsin x$ & $\dfrac{1}{\sqrt{1-x^2}}$ & $\arctan x$ & $\dfrac{1}{1+x^2}$ \\
$\arccos x$ & $-\dfrac{1}{\sqrt{1-x^2}}$ & $\operatorname{arccot}x$ & $-\dfrac{1}{1+x^2}$ \\
\bottomrule
\end{tabular}
\end{center}
\noindent 熟记这张表是求一切导数的基础；配合四则法则、链式法则、反函数法则，就可求初等函数在定义区间内的任意阶导数。

\section{高阶导数}

\subsection{定义与记法}
函数 $y=f(x)$ 的导数 $f'(x)$ 仍是 $x$ 的函数。若 $f'(x)$ 在点 $x$ 处可导，则称 $f'(x)$ 的导数为 $f(x)$ 的\textbf{二阶导数}，记作
\[
y'',\quad f''(x),\quad \frac{\mathrm{d}^2y}{\mathrm{d}x^2};
\]
类似地定义三阶、四阶导数，$n$ 阶导数记作 $y^{(n)}$、$f^{(n)}(x)$、$\dfrac{\mathrm{d}^ny}{\mathrm{d}x^n}$。二阶及二阶以上的导数统称\textbf{高阶导数}。

\subsection{几个常用的 $n$ 阶导数公式}
\keyformula{\begin{aligned}
(e^{ax})^{(n)}&=a^{n}e^{ax}, &\quad (\sin x)^{(n)}&=\sin\Bigl(x+\frac{n\pi}{2}\Bigr),\\
(\cos x)^{(n)}&=\cos\Bigl(x+\frac{n\pi}{2}\Bigr), &\quad (x^{\mu})^{(n)}&=\mu(\mu-1)\cdots(\mu-n+1)x^{\mu-n},\\
(\ln x)^{(n)}&=\frac{(-1)^{n-1}(n-1)!}{x^{n}}. & &
\end{aligned}}

\corepoint{求 $n$ 阶导数的一般方法是“算前几阶、找规律、用数学归纳法证明”。$\sin x,\cos x$ 每求导一次相位前进 $\frac\pi2$，这是周期性函数的高阶导数的常用技巧；$n$ 从 $1$ 开始时 $\sin^{(n)}$ 的公式同样成立。}

\subsection{莱布尼茨公式}
若 $u=u(x)$、$v=v(x)$ 均有 $n$ 阶导数，则
\keyformula{(uv)^{(n)}=\sum_{k=0}^{n}\mathrm{C}_n^{k}\,u^{(n-k)}\,v^{(k)}}
\noindent 其中 $u^{(0)}=u$，$v^{(0)}=v$。形式与二项式展开定理完全相同。

\corepoint{莱布尼茨公式适用于“两个因子之积、其中一个因子（如 $x^m$、多项式）求若干阶导后变零”的情形：多项式 $x^m$ 的 $n$ 阶导在 $n>m$ 时为零，所以和式只剩前面有限项。用错方向的典型是把 $e^x\ln(1+x)$ 硬套莱布尼茨——那不划算。}

\section{隐函数及由参数方程所确定的函数的导数·相关变化率}

\subsection{隐函数的导数}
若方程 $F(x,y)=0$ 确定了 $y$ 是 $x$ 的函数，则称 $y$ 为 $x$ 的\textbf{隐函数}。求导方法：把 $y$ 看作 $x$ 的函数（$y=y(x)$），对方程两边同时关于 $x$ 求导，遇到 $y$ 时按链式法则乘以 $y'$，得到含 $y'$ 的方程，解出 $y'$。

\corepoint{隐函数求导的三个要点：(1) 记住 $y$ 是 $x$ 的函数，$\frac{\mathrm{d}}{\mathrm{d}x}[f(y)]=f'(y)y'$，如 $\frac{\mathrm{d}}{\mathrm{d}x}(y^2)=2yy'$；(2) 结果是含 $x,y$ 的式子，通常保留 $y$（不必解出 $y$）；(3) 求二阶导时，把一阶导的结果\textbf{代入}后再求导，不要停下。求曲线 $F(x,y)=0$ 在点 $(x_0,y_0)$ 处的切线，可用 $y'|_{(x_0,y_0)}$ 直接代入。}

\subsubsection{对数求导法}
对 $y=u(x)^{v(x)}$（$u>0$）或幂连乘、连除式，先在两边取对数 $\ln y=v\ln u$，再两边求导：
\[
\frac{y'}{y}=v'\ln u+v\frac{u'}{u}\ \Longrightarrow\ y'=u^{v}\Bigl(v'\ln u+\frac{vu'}{u}\Bigr).
\]

\corepoint{对数求导法的适用信号：幂指函数（底与指数都含 $x$）、多个因子连乘连除、含根式与乘方。它把乘除化为加减、把指数化为乘数。结果是 $y'$ 的表达式里要\textbf{回乘 $y$ 本身}，别漏掉这个因子（忘记回代是最常见错误）。}

\subsection{由参数方程所确定的函数的导数}
设 $x=\varphi(t)$、$y=\psi(t)$ 都可导且 $\varphi'(t)\neq0$，则
\keyformula{\frac{\mathrm{d}y}{\mathrm{d}x}=\frac{\frac{\mathrm{d}y}{\mathrm{d}t}}{\frac{\mathrm{d}x}{\mathrm{d}t}}=\frac{\psi'(t)}{\varphi'(t)}}
\noindent 二阶导数（对 $x$ 再求一次导，注意 $x$ 仍是 $t$ 的函数）：
\keyformula{\frac{\mathrm{d}^2y}{\mathrm{d}x^2}=\frac{\mathrm{d}}{\mathrm{d}t}\Bigl(\frac{\mathrm{d}y}{\mathrm{d}x}\Bigr)\cdot\frac{1}{\varphi'(t)}
=\frac{\psi''(t)\varphi'(t)-\psi'(t)\varphi''(t)}{\bigl[\varphi'(t)\bigr]^{3}}}

\corepoint{二阶导数\textbf{绝不能}写成 $\dfrac{\psi''(t)}{\varphi''(t)}$。正确做法：先把一阶导 $\dfrac{\mathrm{d}y}{\mathrm{d}x}$ 看成 $t$ 的函数，对它关于 $t$ 求导，再乘以 $\dfrac{\mathrm{d}t}{\mathrm{d}x}=\dfrac{1}{\varphi'(t)}$。这一条几乎每份卷子都在考。}

\subsection{相关变化率}
设 $x=x(t)$、$y=y(t)$ 都可导，而 $x$ 与 $y$ 之间存在关系 $F(x,y)=0$。对 $t$ 求导得
\[
\frac{\mathrm{d}F}{\mathrm{d}x}\cdot\frac{\mathrm{d}x}{\mathrm{d}t}+\frac{\mathrm{d}F}{\mathrm{d}y}\cdot\frac{\mathrm{d}y}{\mathrm{d}t}=0,
\]
从而可由一个变化率求出另一个变化率。这类问题称为\textbf{相关变化率}问题。

\corepoint{解相关变化率问题的步骤：(1) 建立变量之间的方程（几何关系、物理定律）；(2) 两边同时对 $t$ 求导（所有变量都是 $t$ 的函数）；(3) 代入某一时刻的已知值求出待求变化率。注意“代入数值”必须放在求导\textbf{之后}——先代入会把变量变成常数而使导数为 0。}

\section{函数的微分}

\subsection{微分的定义}
\textbf{定义}\quad 设函数 $y=f(x)$ 在某区间内有定义，$x_0$ 及 $x_0+\Delta x$ 都在该区间内。如果函数的增量可表示为
\[
\Delta y=f(x_0+\Delta x)-f(x_0)=A\,\Delta x+o(\Delta x),
\]
其中 $A$ 是不依赖于 $\Delta x$ 的常数，则称 $f(x)$ 在点 $x_0$ 处\textbf{可微}，而 $A\,\Delta x$ 称为 $f(x)$ 在点 $x_0$ 处的\textbf{微分}，记作 $\mathrm{d}y$，即 $\mathrm{d}y=A\,\Delta x$。

\keyformula{f(x)\ \text{在}\ x_0\ \text{处可微}\iff f(x)\ \text{在}\ x_0\ \text{处可导，且}\ \mathrm{d}y=f'(x_0)\,\Delta x}

\noindent 又因 $\mathrm{d}x=\Delta x$（自变量的微分等于自变量的增量），所以
\keyformula{\mathrm{d}y=f'(x)\,\mathrm{d}x,\qquad \frac{\mathrm{d}y}{\mathrm{d}x}=f'(x)}
\noindent 后一式说明：\textbf{导数就是函数的微分与自变量的微分之商}，这就是“微商”一词的来历。

\corepoint{一元函数“可微 $\iff$ 可导”，二者是同一件事；但 $\Delta y$ 与 $\mathrm{d}y$ 不同：$\Delta y$ 是真正的函数增量，$\mathrm{d}y=f'(x)\Delta x$ 是它的\textbf{线性主部}，$\Delta y=\mathrm{d}y+o(\Delta x)$。$\mathrm{d}y$ 依赖于 $x$ 与 $\Delta x$ 两个量，$\mathrm{d}x$ 与 $\Delta x$ 恒等。求 $\mathrm{d}y$ 时不要忘记最后的 $\mathrm{d}x$。}

\subsection{微分的几何意义}
在曲线 $y=f(x)$ 上取点 $M(x_0,y_0)$，当横坐标增加 $\Delta x$ 时，曲线上的纵坐标增量为 $\Delta y=NT$，而切线上相应的纵坐标增量为 $\mathrm{d}y$。因此 $\mathrm{d}y$ 是曲线在点 $M$ 处\textbf{切线的纵坐标增量}，即“以直代曲”的线性近似。

\subsection{基本初等函数的微分公式与微分运算法则}
把求导公式两边同乘 $\mathrm{d}x$ 即得微分公式，如
\[
\mathrm{d}(x^{\mu})=\mu x^{\mu-1}\mathrm{d}x,\quad \mathrm{d}(\sin x)=\cos x\,\mathrm{d}x,\quad
\mathrm{d}(\arctan x)=\frac{\mathrm{d}x}{1+x^2},\quad \mathrm{d}(\ln x)=\frac{\mathrm{d}x}{x}.
\]
微分运算法则：
\keyformula{\mathrm{d}(u\pm v)=\mathrm{d}u\pm\mathrm{d}v,\qquad \mathrm{d}(uv)=v\,\mathrm{d}u+u\,\mathrm{d}v,\qquad \mathrm{d}\Bigl(\frac uv\Bigr)=\frac{v\,\mathrm{d}u-u\,\mathrm{d}v}{v^{2}}}

\subsubsection{复合函数的微分与一阶微分形式不变性}
设 $y=f(u)$ 可微，$u=g(x)$ 可微，则 $\mathrm{d}y=f'(u)\,\mathrm{d}u$。注意这里 $u$ 既可以是自变量，也可以是另一个变量的函数，形式都是 $f'(u)\mathrm{d}u$，这就是\textbf{一阶微分形式不变性}。

\corepoint{一阶微分形式不变性指的是“形式”不变（$\mathrm{d}y=f'(u)\mathrm{d}u$ 恒成立），不是说导数不变（$y'_x=y'_u u'_x$ 差一个 $u'_x$）。它使凑微分（如 $\cos2x\,\mathrm{d}x=\frac12\mathrm{d}(\sin2x)$）成为可能，是第 4 章换元积分法的基础。}

\subsection{微分在近似计算中的应用}
当 $|\Delta x|$ 很小时，
\keyformula{\Delta y\approx \mathrm{d}y=f'(x_0)\Delta x,\qquad f(x_0+\Delta x)\approx f(x_0)+f'(x_0)\Delta x}
\noindent 特别地，取 $x_0=0$ 得 $f(x)\approx f(0)+f'(0)x$，由此得到常用近似式：
\keyformula{(1+x)^{\alpha}\approx1+\alpha x,\quad \sin x\approx x,\quad \tan x\approx x,\quad e^{x}\approx1+x,\quad \ln(1+x)\approx x\quad(|x|\ \text{很小})}

\noindent 误差估计：绝对误差 $|\Delta y|\approx|\mathrm{d}y|=|f'(x)|\,\Delta x$，相对误差 $\dfrac{|\mathrm{d}y|}{|y|}$。

\corepoint{近似公式的前提是 $|x|$（或 $|\Delta x|$）足够小，$\alpha$ 可为任意实数；$x$ 用弧度。用微分作近似时，$\Delta x$ 要取成“方便计算的小量”（如由 $\sqrt{4.02}$ 想到 $x_0=4$、$\Delta x=0.02$）。误差估计的常见坑：多个量相乘时相对误差要相加，不能只算一项。}
